%!TEX root = ../manuscript.tex

\begin{lemma}[Deletion-Monotonicity]\label{lem:del-mon}
Let $w$ be a Gauss word, and let $\mu, \nu$ be two distinct
bigon moves available at $w$ whose chord sets are disjoint.
Then $\nu$ is available at $\mu(w)$, $\mu$ is available at
$\nu(w)$, and $\mu(\nu(w))=\nu(\mu(w))$.
In other words, disjoint bigon moves commute.
\end{lemma}

\begin{proof}
The key observation used throughout: the validity of a bigon move at a
Gauss word is determined by the positions and signs of the endpoints
of its own chord set relative to all other endpoints. Specifically,
the R1-down condition requires two designated endpoints to be
cyclically adjacent, and the bigon R2-down condition requires two
designated arcs to contain no other chord endpoints. In both cases,
\emph{deleting endpoints of chords outside the move's chord set can
only remove obstructions, never create them}. We make this precise in
three cases.

\paragraph{Case 1: R1-down $\circ$ R1-down.}
Let $\mu$ delete the adjacent opposite-sign pair on chord $a$ and
$\nu$ delete the analogous pair on chord $c$, with $a\neq c$.
The two tokens of $c$ are cyclically adjacent at $w$ (for $\nu$ to be
available). Removing the two tokens of $a$ from the cyclic word
preserves the adjacency of $c$'s tokens: if $c$'s tokens were adjacent
at positions $(p, p+1)$ in $w$, then after removing $a$'s two tokens
(which occupy a disjoint pair of positions), the two tokens of $c$
remain cyclically adjacent in the contracted word. Thus $\nu$ is
available at $\mu(w)$. The symmetric argument holds with $\mu, \nu$
reversed. Disjoint R1-downs trivially commute to the same word.

\paragraph{Case 2: R1-down $\circ$ R2-down.}
Let $\mu = \text{R2-down}(a, b)$ with $a$ directly nesting $b$ at $w$,
and $\nu = \text{R1-down}(c)$ with $c \notin \{a, b\}$.
For $\mu$'s availability at $\nu(w)$: the bigon condition for $(a, b)$
at $w$ states that the two outer arcs of $(a, b)$ (the cyclic intervals
$(a_1, b_1)$ and $(b_2, a_2)$) contain no endpoints of any chord other
than $a, b$. In particular, these arcs contain no endpoint of $c$.
Removing $c$'s two endpoints (which lie elsewhere) does not introduce
any new endpoints into the outer arcs, so the bigon condition for
$(a, b)$ continues to hold in $\nu(w)$. Hence $\mu$ is available at
$\nu(w)$.
For $\nu$'s availability at $\mu(w)$: the two tokens of $c$ are
cyclically adjacent in $w$. They are cyclically adjacent via an
interval that contains no endpoint of any other chord (since they are
adjacent). Removing $a$ and $b$ therefore does not affect the
adjacency of $c$'s tokens (the interval between them remains empty).
Hence $\nu$ is available at $\mu(w)$.
The two orders of application produce the same deletion of
$\{a, b, c\}$ tokens and therefore commute.

\paragraph{Case 3: R2-down $\circ$ R2-down.}
Let $\mu = \text{R2-down}(a, b)$ and $\nu = \text{R2-down}(c, d)$ with
$\{a, b\} \cap \{c, d\} = \emptyset$.
The bigon condition for $(a, b)$ at $w$ states that the two outer
arcs of $(a, b)$ are empty of \emph{all} chord endpoints other than
$a, b$'s own; in particular, empty of endpoints of $c$ and $d$.
Removing the four endpoints of $c$ and $d$ (which lie elsewhere) does
not introduce new endpoints into the outer arcs of $(a, b)$, so the
bigon condition for $(a, b)$ persists in $\nu(w)$, and $\mu$ is
available at $\nu(w)$. The symmetric argument shows $\nu$ is available
at $\mu(w)$. Both orderings remove the same four chords $\{a, b, c, d\}$
and leave the other chords with unchanged relative positions; hence the
two resulting words are identical.

Every pair of disjoint bigon moves falls into one of the three cases
above, completing the proof.
\end{proof}
